\section{Формулирование математической постановки задачи}

Аномалиями в поведении пользователей Kubernetes-кластера можно назвать действия, имеющие нетипичные признаки и (или) нетипичные паттерны взаимодействия с ресурсами кластера. Современные кластеры, развёрнутые на основе платформы Kubernetes, регистрируют все обращения к кластеру через механизм audit-логирования, реализованный в компоненте kube-apiserver \cite{KubernetesAudit}.

Задача обнаружения аномалий в audit-логах Kubernetes формализуется как задача обнаружения отклонений. При наличии размеченной тестовой выборки она может рассматриваться как задача бинарной классификации. События рассматриваются как принадлежащие одному из двух классов: нормальные и аномальные. Введём формальные статистические гипотезы: 

\begin{itemize}
    \item Нулевая гипотеза $H_0$ (по умолчанию): событие является нормальным, т.е. $y=0$.
    \item Альтернативная гипотеза $H_1$: событие является аномальным, т.е. $y=1$.
\end{itemize}

В рамках этих гипотез события можно классифицировать следующим образом:

\begin{itemize}
  \item $\mathrm{TP}$ (true positives) --- корректно обнаруженные аномалии (атаку);
  \item $\mathrm{FP}$ (false positives) --- ложные тревоги (нормальные события, помеченные как аномальные);
  \item $\mathrm{FN}$ (false negatives) --- пропущенные аномалии (атаки);
  \item $\mathrm{TN}$ (true negatives) --- корректно не помеченные нормальные события.
\end{itemize}

Задача бинарной классификации заключается в определении принадлежности объектов к одному из двух классов и может решаться методами МО с учителем при наличии размеченной выборки \cite{DeepLearningBook}. Обучение с учителем предполагает наличие набора данных, в котором для каждого объекта заранее известна его истинная метка класса. На основе данной выборки строится модель, параметры которой подбираются таким образом, чтобы обеспечить корректную классификацию новых, ранее не наблюдавшихся объектов тестовой выборки.

Однако в задачах информационной безопасности, включая анализ audit-логов Kubernetes, размеченные данные об аномалиях часто ограничены или отсутствуют. В этом случае задача может формулироваться как задача обнаружения отклонений с использованием методов обучения без учителя \cite{AD_Survey}. При таком подходе модель обучается на данных, характеризующих нормальное поведение, а аномалии выявляются как объекты, существенно отличающиеся от изученного распределения признаков.

Рассматриваемая задача выбора оптимальной модели обнаружения аномалий включает как задачу обучения параметров нейронной сети, так и задачу выбора её архитектуры. Эти задачи имеют различные входные данные и критерии оптимальности, поэтому математическая постановка разбивается на две частные подзадачи.

\begin{enumerate}
    \item Первая подзадача заключается в обучении параметров нейронной сети при фиксированной архитектуре. В рамках данной подзадачи необходимо определить значения параметров модели, которые наилучшим образом описывают нормальное поведение пользователей кластера на основе обучающей выборки. Эта задача является задачей оптимизации функции потерь модели.
    \item Вторая подзадача заключается в выборе оптимальной структуры нейронной сети из заданного множества допустимых архитектур. Для каждой архитектуры выполняется обучение параметров, после чего производится оценка качества обнаружения аномалий на тестовой выборке. Таким образом, данная подзадача представляет собой задачу оптимизации по архитектуре модели.
\end{enumerate}

Разделение задачи на две подзадачи позволяет отделить процесс обучения параметров модели от процесса выбора её структуры и обеспечить корректное сравнение различных архитектур по единому критерию качества.

\subsection{Исходные данные}

Каждая запись audit-лога представляет собой структурированное сообщение в формате JSON, содержащее сведения о выполненном запросе, субъекте, инициировавшем операцию, целевом ресурсе и результате выполнения запроса. Audit-лог можно представить в виде упорядоченной во времени последовательности событий, как показано в формуле~\eqref{eq:audit_log_sequence}.

\begin{equation}\label{eq:audit_log_sequence}
    \mathcal{L} = ( e_1, e_2, \dots, e_T )
\end{equation}

В ходе анализа структуры audit-логов были выделены поля, наиболее информативные для последующего применения методов обнаружения аномалий, характеризующие событие $e_t$, как показано в формуле~\eqref{eq:event_structure}. Аналогичный набор признаков использовался в работе \cite{K8NTEXT} для определения контекста между событиями с помощью НС.

\begin{equation}\label{eq:event_structure}
    e_t = (v_t, u_t, ua_t, r_t, sr_t, ns_t, c_t, ip_t, t_t),
\end{equation}

где:
\begin{itemize}
    \item $v_t$ — значение поля \texttt{verb}, определяющее тип операции, выполняемой пользователем или сервисным аккаунтом над ресурсом Kubernetes (например, \texttt{get}, \texttt{list}, \texttt{create}, \texttt{update}, \texttt{delete}, \texttt{watch});
    \item $u_t$ — значение поля \texttt{user.username}, содержащее имя пользователя или сервисного аккаунта, от имени которого был выполнен запрос к API-серверу Kubernetes;
    \item $ua_t$ — значение поля \texttt{userAgent}, содержащее строку идентификации клиента, осуществившего запрос к API (например, \texttt{kubectl}, клиентские библиотеки Kubernetes или автоматизированные сервисы);
    \item $r_t$ — значение поля \texttt{objectRef.resource}, определяющее тип ресурса Kubernetes, к которому осуществляется обращение (например, \texttt{pods}, \texttt{deployments}, \texttt{configmaps}, \texttt{secrets});
    \item $sr_t$ — значение поля \texttt{objectRef.subresource}, определяющее подресурс объекта Kubernetes, если обращение производится не к основному ресурсу, а к его части (например, \texttt{status}, \texttt{scale}, \texttt{log});
    \item $ns_t$ — значение поля \texttt{objectRef.namespace}, определяющее пространство имён Kubernetes, в котором расположен целевой ресурс;
    \item $c_t$ — значение поля \texttt{responseStatus.code}, содержащее HTTP-код ответа API-сервера, характеризующий результат выполнения запроса (например, 200, 201, 403, 404);
    \item $ip_t$ — значение поля \texttt{sourceIPs}, содержащее список IP-адресов источников запроса, с которых был выполнен доступ к API-серверу Kubernetes;
    \item $t_t$ — значение поля \texttt{requestReceivedTimestamp}, содержащее временную метку получения запроса API-сервером и определяющее момент времени, в который событие было зафиксировано в audit-логе.
\end{itemize}

Выбор указанных полей обусловлен тем, что они отражают ключевые характеристики взаимодействия пользователей и сервисов с ресурсами кластера: тип выполняемой операции, субъект действия, используемый клиент, объект доступа и результат выполнения запроса. Совместное использование указанных признаков позволяет формировать последовательности событий, характеризующие типичное поведение субъектов системы, что является основой для последующего применения методов обнаружения аномалий.

Для кодирования событий в числовой вектор используется функция $\Phi$, определяемой формулой~\eqref{encoding_function}, которая преобразует сырую последовательность структурированных записей audit-лога $\mathcal{L} = \{e_1, \dots, e_T\}$ к единому признаковому представлению.

\begin{equation}\label{encoding_function}
    \Phi : \mathcal{E} \rightarrow \mathbb{R}^d,
\end{equation}

где $\mathcal{E}$ — множество всех возможных событий audit-лога. В результате формируется последовательность векторов признаков для каждого события $e_t$, определяемая формулой~\eqref{feature_vector}.

\begin{equation}\label{feature_vector}
    x_t = \Phi(e_t) \in \mathbb{R}^d.
\end{equation}

Таким образом, формируется выборка, определяемая формулой~(\ref{sample}).

\begin{equation}\label{sample}
    \mathcal{X} = \{ x_1, x_2, \dots, x_T \}, \quad x_t \in \mathbb{R}^d.
\end{equation}

Предполагается, что исходная выборка разбивается на три подмножества:

\begin{itemize}
    \item обучающая выборка $\mathcal{X}_{train} \subset \mathcal{X}$, содержащая объекты, соответствующие нормальному поведению пользователей в кластере, используемая для обучения параметров моделей;
    \item валидационная выборка $\mathcal{X}_{val} \subset \mathcal{X}$, используемая для оценки значения функции потерь и выбора оптимальных параметров модели;
    \item тестовая выборка $\mathcal{X}_{test} \subset \mathcal{X}$, содержащая размеченные объекты и используемая для итоговой оценки качества моделей.
\end{itemize}

Пусть задано множество допустимых архитектур нейронных сетей

$$\mathcal{A} = \{ A_\alpha \},$$

Для каждой архитектуры $A_\alpha$ определяется множество параметров модели $\theta$. Модель обозначим как

$$f_{\alpha,\theta} : \mathbb{R}^d \rightarrow \mathbb{R},$$

где

$\alpha$ - структура сети,

$\theta$ - обучаемые параметры.

\subsection{Выбор параметров нейронной сети}

Входными данными подзадачи являются:

\begin{itemize}
    \item обучающая выборка событий $\mathcal{X}_{train}$;
    \item валидационная выборка $\mathcal{X}_{val}$;
    \item фиксированная архитектура нейронной сети $A_\alpha$;
    \item функция потерь модели $\mathcal{F}$.
\end{itemize}

Для определения параметров нейронной сети используется функция потерь, зависящая от выбранной архитектуры модели. Обучение параметров осуществляется на обучающей выборке $\mathcal{X}_{train}$, а значение функции потерь вычисляется на валидационной выборке $\mathcal{X}_{val}$, что позволяет оценивать обобщающую способность модели и предотвращать её переобучение. Оптимальные параметры модели определяются согласно формуле~(\ref{optim_params}).


\begin{equation}\label{optim_params}
    \theta^*(\alpha) = \arg\min_{\theta} \frac{1}{|\mathcal{X}_{val}|} \sum_{x \in \mathcal{X}_{val}} \mathcal{F}(x; \alpha, \theta),
\end{equation}

где $\mathcal{F}$ — функция потерь, зависящая от выбранного типа модели (например, ошибка реконструкции Autoencoder \cite{AE}).

Результатом решения подзадачи является оптимальный набор параметров модели $\theta^*(\alpha)$.

\subsection{Выбор структуры нейронной сети}

Входными данными подзадачи являются:

\begin{itemize}
    \item множество допустимых архитектур нейронных сетей $\mathcal{A}$;
    \item обученные модели $f_{\alpha,\theta^*(\alpha)}$ для каждой архитектуры;
    \item тестовая выборка $\mathcal{X}_{test}$;
    \item вектор меток $\mathcal{Y}_{test}$.
\end{itemize}

Для оценки качества моделей обнаружения аномалий используется метрика площади под PR-кривой (AUC-PR). Важность выбора PR-метрики при дисбалансе классов показана в работе \cite{PRCurve}. PR-кривая отражает зависимость точности (Precision), вычисляемой по формуле~(\ref{precision}), от полноты (Recall), вычисляемой по формуле~(\ref{recall}), при варьировании порога обнаружения.

\begin{equation}\label{precision}
    \mathrm{Precision} = \frac{TP}{TP + FP}.
\end{equation}

\begin{equation}\label{recall}
    \mathrm{Recall} = \frac{TP}{TP + FN}.
\end{equation}


Пусть имеется тестовая выборка $\mathcal{X}_{test}$ с известными метками
$\mathcal{Y}_{test} = \{y_i\}$, $y_i \in \{0,1\}$.

Для каждой архитектуры $A_\alpha$ вычисляется набор показателей аномальности, определяемый формулой~(\ref{indicators_anomaly}).

\begin{equation}\label{indicators_anomaly}
    S_\alpha = \{ s_\alpha(x_1), \dots, s_\alpha(x_n) \},
\end{equation}

где $s_{\alpha}(x_i) = f_{\alpha,\theta^*(\alpha)}(x_i)$ - показатель аномальности события

Качество детекции оценивается с использованием метрики PR-AUC, определяемой формулой~(\ref{pr_auc}).

\begin{equation}\label{pr_auc}
   Q(\alpha) = \mathrm{AUC}_{PR}(\mathcal{Y}_{test}, S_\alpha).
\end{equation}

Тогда задача выбора структуры нейронной сети формализуется формулой~(\ref{optim_arch}).

\begin{equation}\label{optim_arch}
   \alpha^* =\arg\max_{\alpha \in \mathcal{A}}\mathrm{AUC}_{PR}(\mathcal{Y}_{test}, S_\alpha).
\end{equation}

\subsection{Обобщённая формулировка задачи}

Таким образом, задача сводится к двум подзадачам:

\begin{enumerate}
    \item выбор параметров сети при фиксированной архитектуре на наборе $\mathcal{X}_{train}$;
    \item выбор оптимальной структуры сети по критерию максимизации PR-AUC на наборе $\mathcal{X}_{test}$.
\end{enumerate}


Итоговая оптимизационная задача задаётся системой формул~(\ref{math_statement_1})–(\ref{math_statement_2}). Оптимальные параметры, полученные из формулы~(\ref{math_statement_1}), используются при решении задачи, представленной формулой~(\ref{math_statement_2}).

\begin{equation}\label{math_statement_1}
   \theta^*(\alpha) = \arg\min_{\theta} \frac{1}{|\mathcal{X}_{val}|} \sum_{x \in \mathcal{X}_{val}} \mathcal{F}(x; \alpha, \theta),
\end{equation}

\begin{equation}\label{math_statement_2}
   \alpha^* =\arg\max_{\alpha \in \mathcal{A}}\mathrm{AUC}_{PR}(\mathcal{Y}_{test}, S_\alpha).
\end{equation}

Полученная архитектура $A_{\alpha^*}$ и параметры $\theta^*(\alpha^*)$ задают оптимальную модель обнаружения аномалий в рамках выбранного множества архитектур $\mathcal{A}$.

\subsection{Методы решения задачи}

Решение поставленной задачи предполагает последовательное выполнение нескольких этапов: подготовку и предобработку исходных данных, формирование признакового описания событий, обучение моделей нейронных сетей и оценку их качества.

На этапе подготовки данных формируется набор событий audit-логов $\mathcal{X}_{train}$ и $\mathcal{X}_{test}$. Для тестового набора выполняется разметка событий на классы (нормальные и аномальные), при этом набор должен включать как типичные операции легитимных пользователей, так и возможные сценарии атак, попыток эскалации привилегий, массового сканирования ресурсов и несанкционированного доступа.

На этапе предобработки данных осуществляется преобразование исходных событий audit-логов в численные векторы признаков, пригодные для обучения НС. Основные задачи включают:

\begin{itemize}
    \item \textbf{Очистка и унификация данных:} устранение пропусков, нормализация числовых значений, приведение категориальных полей к единому формату. Это обеспечивает согласованность представления событий и снижает вероятность ошибок при обучении моделей.

    \item \textbf{Кодирование категориальных признаков:} преобразование различных типов категориальных данных (например, идентификаторы пользователей, типы операций, виды ресурсов) в численные представления. Возможны подходы, такие как one-hot кодирование, хеширование категорий или обучаемые embedding-слои. Это позволяет модели выявлять закономерности в действиях пользователей и взаимодействиях с ресурсами.

    \item \textbf{Построение временных и контекстных признаков:} выделение информации о временных паттернах событий (циклические характеристики, рабочее/нерабочее время), истории действий пользователя, последовательности обращений к ресурсам и агрегированных статистик по различным объектам кластера. Такие признаки повышают информативность представления событий и улучшают способность модели обнаруживать нетипичные последовательности действий.

    \item \textbf{Интеграция флагов и дополнительных индикаторов:} формирование булевых или численных индикаторов, отражающих специфические характеристики запросов, состояния ресурсов или контекста операции. Например, можно учитывать доступ к чувствительным ресурсам, статус ответа системы или происхождение запроса. Эти признаки помогают выявлять потенциально аномальные сценарии работы.

\end{itemize}

В результате предобработки каждое событие представляется в виде вектора признаков фиксированной размерности, содержащего числовые, категориальные и контекстные компоненты.

Для классификации событий целесообразно использовать НС различных архитектур, включая RNN \cite{RNN}, LSTM \cite{LSTM} и автокодировщики \cite{AE}. Для каждой архитектуры сети $A_\alpha$ выполняется обучение на наборе $\mathcal{X}_{train}$ с целью минимизации функции потерь $\mathcal{F}$, а затем рассчитывается показатель аномальности $s\alpha(x)$ для тестовых событий $\mathcal{X}_{test}$. Качество моделей оценивается с использованием метрики PR-AUC, что особенно важно при дисбалансе классов. По результатам сравнительного анализа выбирается оптимальная архитектура $\alpha^*$ и соответствующие параметры $\theta^(\alpha^*)$, формирующие рабочую модель обнаружения аномалий в Kubernetes-кластере.